\documentclass[
  reprint,
  amsmath,
  amssymb,
  aps,
  pra,
  nofootinbib,
  superscriptaddress,
  floatfix
]{revtex4-2}

\usepackage[T1]{fontenc}
\usepackage{lmodern}
\usepackage{graphicx}
\usepackage{booktabs}
\usepackage{bm}
\usepackage{mathtools}
\usepackage{xcolor}
\usepackage{hyperref}
\hypersetup{hidelinks}

\newcommand{\placeholder}[1]{\textcolor{red!70!black}{\ifmmode\left[#1\right]\else\textsf{[#1]}\fi}}
\newcommand{\tentative}[1]{\textcolor{red!70!black}{#1}}
\newcommand{\Tr}{\operatorname{Tr}}
\newcommand{\dd}{\mathrm{d}}
\newcommand{\ii}{\mathrm{i}}
\newcommand{\GKBA}{\mathrm{GKBA}}
\newcommand{\EOM}{\mathrm{EOM}}
\newcommand{\Reop}{\operatorname{Re}}
\newcommand{\argmin}{\operatorname*{arg\,min}}
\newcommand{\ket}[1]{\left|#1\right\rangle}

\begin{document}

\title{Regularized Equal-Time Electron--Phonon Dynamics for Quantum Simulation}

\author{Jake Strobel}
\author{Jacopo Simoni}
\author{Gabriele Riva}
\author{Yuan Ping}
\affiliation{University of Wisconsin--Madison}

\date{\today}

\begin{abstract}
We study whether the equal-time non-Markovian electron--phonon equations of motion introduced by Riva, Simoni, and Ping can be turned into a quantum-simulation target in the high-interaction regime. The working object is the coupled ODE system for the electronic one-body density matrix, phonon covariance, anomalous phonon covariance, coherent phonon amplitude, and electron--phonon correlation. This planning template fixes the source equations and records candidate solver routes, separated into near-term variational approaches and fault-tolerant approaches. No stability, complexity, or quantum-advantage result is claimed in this draft.
\end{abstract}

\maketitle

\section{Objective}

The target problem is the direct numerical solution of the equal-time electron--phonon ODEs on page 3 of Ref.~\cite{RivaSimoniPing2026}, followed by a controlled high-\(U\) regularization study. The first planning deliverable is a list of plausible classical, NISQ, and fault-tolerant routes, together with the diagnostics needed to decide which route is technically defensible for the tested \(L=2\) case.

\section{Target ODE system}
\label{sec:target_ode}

The dynamical variables are the electronic one-body density matrix \(\rho_{12}\), the phonon covariance \(\delta\rho_{\mathbf q\mathbf q'}\), the anomalous phonon covariance \(\delta\bar\rho_{\mathbf q'\mathbf q}\), the electron--phonon correlation \(\delta\rho^{\mathbf q}_{12}\), and the coherent phonon amplitude \(B_{\mathbf q}\). We take the source system to be Eq.~(14a)--(14e), with the effective Hamiltonian in Eq.~(15), of Ref.~\cite{RivaSimoniPing2026}. The working algorithmic split is that Eqs.~\eqref{eq:paper_v_rho_electron}--\eqref{eq:paper_v_delta_rho12q} define the quantum-computer target, while Eq.~\eqref{eq:paper_v_B} is a classical coherent-field update coupled back through \(\tilde h(t)\).

\begin{widetext}
\begin{subequations}
\label{eq:paper_v_source_eom}
\begin{align}
\dot{\rho}_{12}
&=
-\ii[\tilde h,\rho]_{12}
+\left[
-\ii\sum_{\mathbf q}\sum_3
\left(
g_{13}^{\mathbf q-}\delta\rho_{32}^{\mathbf q}
+g_{31}^{\mathbf q+}{\delta\rho_{23}^{\mathbf q}}^{*}
\right)
+\mathrm{H.c.}
\right],
\label{eq:paper_v_rho_electron}
\\
\delta\dot{\rho}_{\mathbf q\mathbf q'}
&=
-\ii(\omega_{\mathbf q}-\omega_{\mathbf q'})
\delta\rho_{\mathbf q\mathbf q'}
-\ii\sum_{12}
\left(
-g_{12}^{\mathbf q'-}\delta\rho_{21}^{\mathbf q}
+g_{12}^{\mathbf q+}{\delta\rho_{21}^{\mathbf q'}}^{*}
\right),
\label{eq:paper_v_delta_rho}
\\
\delta\dot{\bar\rho}_{\mathbf q'\mathbf q}
&=
-\ii(\omega_{\mathbf q'}+\omega_{\mathbf q})
\delta\bar\rho_{\mathbf q'\mathbf q}
-\ii\sum_{12}
\left(
g_{12}^{\mathbf q+}\delta\rho_{12}^{\mathbf q'}
+g_{12}^{\mathbf q'+}\delta\rho_{12}^{\mathbf q}
\right),
\label{eq:paper_v_delta_rho_bar}
\\
\delta\dot{\rho}_{12}^{\mathbf q}
&=
-\ii\left(
[\tilde h,\delta\rho^{\mathbf q}]_{12}
+\omega_{\mathbf q}\delta\rho_{12}^{\mathbf q}
\right)
-\ii\sum_{\mathbf q'}\sum_{34}
g_{34}^{\mathbf q'+}
(\delta_{\mathbf q\mathbf q'}+\delta\rho_{\mathbf q\mathbf q'})
(\delta_{14}-\rho_{14})\rho_{32}
\nonumber\\
&\quad
+\ii\sum_{\mathbf q'}\sum_{34}
g_{34}^{\mathbf q'+}\delta\rho_{\mathbf q\mathbf q'}
(\delta_{32}-\rho_{32})\rho_{14}
-\ii\sum_{3\mathbf q'}
g_{23}^{\mathbf q'-}\delta\bar\rho_{\mathbf q'\mathbf q}\rho_{31}
+\ii\sum_{3\mathbf q'}
g_{32}^{\mathbf q'-}\delta\bar\rho_{\mathbf q'\mathbf q}\rho_{13},
\label{eq:paper_v_delta_rho12q}
\\
\dot B_{\mathbf q}
&=
-\ii\left(
\omega_{\mathbf q}B_{\mathbf q}
+\sum_{12}g_{12}^{\mathbf q+}\rho_{12}
\right),
\label{eq:paper_v_B}
\\
\tilde h_{12}(t)
&=
\varepsilon_1\delta_{12}
+\sum_{\mathbf q}g_{12}^{\mathbf q-}
\left(B_{\mathbf q}+B_{-\mathbf q}^{*}\right)
+v_{12}^{\mathrm{ext}}(t).
\label{eq:paper_v_htilde}
\end{align}
\end{subequations}
\end{widetext}

\section{Baseline information needed}
\label{sec:baseline_needed}

Let
\begin{equation}
u(t)=\operatorname{vec}
\left(
\rho,\delta\rho,\delta\bar\rho,\delta\rho^{\mathbf q}
\right)
\end{equation}
denote the real-doubled vectorization of the quantum-target sector, Eqs.~\eqref{eq:paper_v_rho_electron}--\eqref{eq:paper_v_delta_rho12q}. The coherent phonon \(B_{\mathbf q}(t)\) is kept as a classical auxiliary field. The coupled system can be viewed as
\begin{equation}
\dot u=F_Q(u,B,t),
\qquad
\dot B=F_B(u,B,t).
\label{eq:paper_v_vector_ode}
\end{equation}
Before choosing a quantum algorithm, the classical \(L=2\) solve should determine the vector dimension of \(u\), sparsity, stiffness, nonlinear tensor contractions, invariant drift, and the precise high-\(U\) failure mode.

\section{NISQ ideas}
\label{sec:nisq_ideas}

\begin{itemize}
\item Variational ODE residual solver: represent \(u(t)\approx r(t)z(\bm\theta(t))\) with a normalized circuit state plus scale and minimize \(\|\dot u-F_Q(u,B,t)\|^2\) over parameter velocities, with \(B(t)\) supplied by a classical update.
\item McLachlan-style projected dynamics: use tangent-space metric and force measurements to update \(\bm\theta(t)\), treating the ODE residual as the projected drift.
\item pVQD-style time-step matching: choose \(\bm\theta_{k+1}\) so the circuit state matches an explicit or implicit Euler/RK step \(u_k+\Delta t\,F_Q(u_k,B_k,t_k)\).
\item Variational linear-solver subroutine: use VQLS-like circuits only inside implicit time steps, Newton corrections, or regularized linearized updates.
\item Tensor-factor variational ansatz: keep separate circuit or tensor factors for electronic density, phonon covariance, anomalous covariance, and electron--phonon correlation sectors; keep the coherent phonon field classical.
\item Hybrid contraction path: compute nonlinear products in \(F_Q(x,B,t)\) classically from measured or reconstructed low-dimensional observables, then use quantum circuits only for the largest linear-algebra subproblem.
\item Measurement-first diagnostic route: use the NISQ circuit only to estimate invariants, residual norms, and selected observables before claiming full trajectory propagation.
\end{itemize}

\section{McLachlan pathway}
\label{sec:mclachlan_pathway}

The direct NISQ propagation route is a McLachlan projection of the
real-doubled ODE sector.  It does not replace the coupled ODE solve by a
closed Hamiltonian simulation.  It chooses a variational manifold
\(\mathcal M=\{x_{\bm\vartheta}\}\subset\mathbb R^N\) for the
quantum-target vector \(x\), evaluates the drift \(F_Q(x_{\bm\vartheta},B,t)\)
at the current classical field \(B(t)\), and projects that drift onto the
tangent space of \(\mathcal M\).  With
\begin{equation}
J_{\bm\vartheta}
=
\left[
\partial_{\vartheta_1}x_{\bm\vartheta}\ \cdots\
\partial_{\vartheta_m}x_{\bm\vartheta}
\right],
\end{equation}
the projected velocity is
\begin{equation}
\dot{\bm\vartheta}
=
\argmin_{\bm\nu}
\left\|J_{\bm\vartheta}\bm\nu
-F_Q(x_{\bm\vartheta},B,t)\right\|_2^2 ,
\label{eq:paper_v_mclachlan_argmin}
\end{equation}
or
\begin{equation}
\left(J_{\bm\vartheta}^{T}J_{\bm\vartheta}+\Lambda\right)
\dot{\bm\vartheta}
=
J_{\bm\vartheta}^{T}F_Q(x_{\bm\vartheta},B,t)
\label{eq:paper_v_mclachlan_normal}
\end{equation}
with the same kind of thresholded or damped inverse used in regularized
variational dynamics.

For this system the useful ansatz should carry separate sector scales, since
\(\rho\), \(\delta\rho\), \(\delta\bar\rho\), and
\(\delta\rho^{\mathbf q}\) need not have comparable norms.  Let
\begin{equation}
x_{\bm\vartheta}
=
\sum_s r_s z_s(\bm\theta),
\qquad
z_s=\Pi_s z(\bm\theta),
\end{equation}
where
\(s\in\{\rho,\delta\rho,\delta\bar\rho,\delta\rho^{\mathbf q}\}\) and
\(\Pi_s\) are orthogonal sector projectors.  Define the sector norm and
horizontal tangent by
\begin{align}
n_s&=\langle z_s,z_s\rangle_{\mathbb R},
\\
\bar T_{i,s}
&=
\Pi_s\partial_i z
-
\frac{
\langle z_s,\Pi_s\partial_i z\rangle_{\mathbb R}
}{n_s}\,z_s ,
\qquad n_s>0 .
\end{align}
For \(F_s=\Pi_sF_Q(x_{\bm\vartheta},B,t)\), the residual is
\begin{equation}
\begin{aligned}
R_Q
&=
\sum_s \dot r_s z_s
+
\sum_i
\left(\sum_s r_s\bar T_{i,s}\right)\dot\theta_i
\\
&\quad
-F_Q(x_{\bm\vartheta},B,t).
\end{aligned}
\label{eq:paper_v_mclachlan_residual}
\end{equation}
The sector-gauge McLachlan equations are
\begin{align}
n_s\dot r_s&=\Phi_s,
\qquad
\Phi_s=\langle z_s,F_s\rangle_{\mathbb R},
\label{eq:paper_v_sector_scale_eq}
\\
\left(
\sum_s r_s^2\bar M^{(s)}+\Lambda
\right)\dot{\bm\theta}
&=
\sum_s r_s\bar{\bm V}^{(s)},
\label{eq:paper_v_sector_theta_eq}
\end{align}
with
\begin{equation}
\bar M_{ij}^{(s)}
=
\langle\bar T_{i,s},\bar T_{j,s}\rangle_{\mathbb R},
\qquad
\bar V_i^{(s)}
=
\langle\bar T_{i,s},F_s\rangle_{\mathbb R}.
\end{equation}
The scalar-scale equations in Sec.~\ref{sec:pedagogical_solver_map} are the
one-sector special case.

This is where the route differs from ordinary Hamiltonian variational
dynamics.  The tangent metric and sector probabilities are natural QPU-facing
objects.  The small normal-equation solve, the \(B_{\mathbf q}\) update, and
the time integrator are classical.  The bottleneck is the force projection
\(\bar V_i^{(s)}\), because with \(B(t)\) classical the target force has the
schematic form
\begin{equation}
F_Q(x,B,t)
=
A(B,t)x+\mathcal Q(x,x;B,t)+\mathcal C(x,x,x;B,t),
\label{eq:paper_v_force_polynomial}
\end{equation}
so
\begin{equation}
V_i
=
r\langle T_i,Az\rangle_{\mathbb R}
+r^2\langle T_i,\mathcal Q(z,z)\rangle_{\mathbb R}
+r^3\langle T_i,\mathcal C(z,z,z)\rangle_{\mathbb R}.
\label{eq:paper_v_force_projection_layers}
\end{equation}
The first term resembles standard variational simulation.  The quadratic and
cubic terms require an explicit access strategy for nonlinear products.

The planning split is therefore four force-evaluation routes:
\begin{enumerate}
\item reconstruct \(x\) for \(L=2\), compute \(F_Q\) classically, and use the
quantum circuit as a variational compression or residual diagnostic;
\item measure only selected sparse or low-rank contractions when the coupling
tensors make the projection \(\langle T_i,F_Q\rangle\) compact;
\item estimate multilinear products with multiple circuit copies, which is
conceptually direct but shot- and noise-sensitive on NISQ hardware;
\item move to a fault-tolerant oracle or block encoding for \(F_Q\) or its
Jacobian \(J_Q=\partial F_Q/\partial x\).
\end{enumerate}
Only the first route is immediately honest for the \(L=2\) dimer.  A scalable
NISQ claim needs a non-tomographic method for the multilinear force projections
in Eq.~\eqref{eq:paper_v_force_projection_layers}.

The minimal \(L=2\) benchmark should initialize \((x_0,B_0)\), fit
\((\bm r_0,\bm\theta_0)\) to \(x_0\), propagate
\((\bm r,\bm\theta,B)\) by Eqs.~\eqref{eq:paper_v_sector_scale_eq} and
\eqref{eq:paper_v_sector_theta_eq}, and compare the reconstructed trajectory
against Radau/BDF/DOP853 classical references.  The required diagnostics are
the ODE residual \(\|R_Q\|\), sector residuals \(\|\Pi_sR_Q\|\), trajectory
errors for \(x\) and \(B\), \(\Tr\rho\), Hermiticity, positivity drift, and
growth of \(\delta\rho^{\mathbf q}\).  The technical validation is projected
variational propagation of the quantum-target sector with classical coherent
feedback, not a quantum-advantage claim.

\section{Fault-tolerant ideas}
\label{sec:ft_ideas}

\begin{itemize}
\item Sparse ODE simulation after linearization: linearize the quantum-target sector of Eq.~\eqref{eq:paper_v_vector_ode} around a trajectory or fixed point and block-encode the resulting Jacobian, with \(B(t)\) treated as a classical control field.
\item Carleman or Koopman lifting: embed the nonlinear polynomial ODE into a larger approximately linear system, then use fault-tolerant linear ODE or matrix-exponential tools.
\item Quantum linear systems for implicit steps: solve \((I-\Delta t\,J_k)\delta y=b_k\) with a fault-tolerant QLSA inside Newton--Krylov or backward-Euler integration.
\item LCU/QSP matrix functions: apply \(e^{A\Delta t}\), resolvents, or regularized propagators when a stable linear generator \(A\) is available.
\item Amplitude estimation for observables: use FT amplitude estimation to reduce sampling cost for conserved quantities, spectra, or residual norms.
\item Full Liouville-space embedding: encode the density/correlation vector from Eqs.~\eqref{eq:paper_v_rho_electron}--\eqref{eq:paper_v_delta_rho12q} directly and propagate with a block-encoded non-Hermitian or dilated generator when the truncation is linear.
\item FT benchmark role: use fault-tolerant algorithms as an asymptotic comparison even if the first implementation remains classical or NISQ.
\end{itemize}

\section{Numerical plan}

\paragraph*{Classical/quantum split.}
Flatten the complex matrices in Eqs.~\eqref{eq:paper_v_rho_electron}--\eqref{eq:paper_v_delta_rho12q} into the quantum-target real state \(u(t)\) by storing real and imaginary parts separately. Propagate \(B_{\mathbf q}\) from Eq.~\eqref{eq:paper_v_B} as a classical auxiliary variable. The first smoke case should preserve the source-paper Holstein-dimer symmetry before adding high-\(U\) stress tests.

\paragraph*{Invariants.}
The minimum diagnostics are \(\Tr\rho\), Hermiticity of \(\rho\), positivity drift, total energy after the external pulse, norm growth of \(\delta\rho^{\mathbf q}_{12}\), and solver-step rejection counts.

\paragraph*{Regularization slot.}
\placeholder{Specify the high-\(U\) failure mode: density-matrix positivity loss, stiffness, correlation blow-up, energy drift, or self-consistency instability.}

\paragraph*{Quantum encoding slot.}
\placeholder{Choose after the classical \(L=2\) diagnostics identify the dominant bottleneck: state dimension, stiffness, nonlinearity, measurement burden, or invariant preservation.}

\section{Immediate tasks}

\begin{enumerate}
\item Build the \(L=2\) finite-dimensional tensor layout for the quantum-target variables in Eqs.~\eqref{eq:paper_v_rho_electron}--\eqref{eq:paper_v_delta_rho12q} and the classical coherent field in Eq.~\eqref{eq:paper_v_B}.
\item Implement the right-hand side with explicit shape, Hermiticity, and index-convention tests.
\item Reproduce a coherent-dynamics baseline by solving Eqs.~\eqref{eq:paper_v_rho_electron} and~\eqref{eq:paper_v_B} with a classical ODE integrator.
\item Add the full non-Markovian system and record the first high-\(U\) instability signature.
\item Decide whether the first quantum route for Eqs.~\eqref{eq:paper_v_rho_electron}--\eqref{eq:paper_v_delta_rho12q} should be NISQ variational, NISQ measurement-assisted, FT linearized, or FT lifted-linear.
\end{enumerate}

\section{Pedagogical solver map}
\label{sec:pedagogical_solver_map}

The first useful simplification is to separate the object being propagated from
the hardware representation used to propagate it.  The quantum-target state is
not a wavefunction.  After real-doubling,
\begin{equation}
x(t)=\mathcal R\,\operatorname{vec}
\left(
\rho,\delta\rho,\delta\bar\rho,\delta\rho^{\mathbf q}
\right)
\in\mathbb R^N ,
\end{equation}
the target dynamics have the form
\begin{equation}
\dot x=F_Q(x,B,t),
\qquad
\dot B=F_B(x,B,t).
\end{equation}
The coherent phonon field \(B(t)\) remains a classical control variable.  It is
updated from Eq.~\eqref{eq:paper_v_B} and re-enters the quantum-target sector
through \(\tilde h(t)\).  A quantum algorithm for this draft is therefore a
method for representing, projecting, linearizing, or estimating the
\(x\)-sector, with \(B(t)\) supplied by a classical feedback loop.

\paragraph*{Classical diagnostic layer.}
Before a quantum route is meaningful, the \(L=2\) system should be solved as a
classical coupled ODE for \((x,B)\).  The classical solve is not the final
algorithmic claim; it identifies the numerical obstruction that the quantum
method would need to address:
\begin{equation}
\left\{
\dim x,\ \operatorname{nnz} J_Q,\ 
\operatorname{spec} J_Q,\ 
\kappa(I-\Delta t\,J_Q),\
\Delta_{\rm phys}(t)
\right\},
\end{equation}
where \(J_Q=\partial F_Q/\partial x\) and \(\Delta_{\rm phys}\) collects trace,
Hermiticity, positivity, covariance-symmetry, energy, and correlation-growth
defects.  If adaptive explicit solvers, stiff implicit solvers, and
Jacobian-aware methods agree at matched tolerances, the first quantum claim
should be about representation or observable estimation rather than speed.

\paragraph*{NISQ residual layer.}
A near-term variational route should be written as a residual projection, not
as unitary Schrodinger evolution.  Let
\begin{equation}
x(t)\approx r(t)z(\bm\theta(t)),
\qquad
\|z(\bm\theta)\|_2=1 .
\end{equation}
The scale \(r(t)\) is required because the ODE state is not normalized.  The
instantaneous defect is
\begin{equation}
R_Q=
\dot r\,z
+r\sum_i\partial_i z\,\dot\theta_i
-F_Q(rz,B,t),
\end{equation}
and the McLachlan-type update chooses \((\dot r,\dot{\bm\theta})\) by minimizing
\(\|R_Q\|_2^2\).  In the common gauge
\(\operatorname{Re}\langle z,\partial_i z\rangle=0\), this gives the schematic
normal equations
\begin{equation}
\dot r=\operatorname{Re}\langle z,F_Q\rangle,
\qquad
r^2 M_{ij}\dot\theta_j
=r V_i ,
\end{equation}
with
\begin{equation}
M_{ij}=\operatorname{Re}\langle\partial_i z|\partial_jz\rangle,
\qquad
V_i=\operatorname{Re}\langle\partial_i z|F_Q(rz,B,t)\rangle .
\end{equation}
This layer is plausible only if the metric, force vector, nonlinear products,
and selected observables can be measured without reconstructing all of
\(x(t)\).  The strongest near-term target is therefore residual and invariant
estimation for selected observables, with full trajectory reconstruction left
as a small-system diagnostic.

\paragraph*{Fault-tolerant linearization layer.}
Fault-tolerant ODE algorithms become natural after the nonlinear target is
turned into a linear algebra problem.  Around a classical or variational
trajectory \(\bar x(t)\),
\begin{equation}
x(t)=\bar x(t)+\xi(t)
\end{equation}
gives
\begin{equation}
\dot\xi
=J_Q(\bar x(t),B(t),t)\xi+s(t)+O(\|\xi\|^2),
\end{equation}
where \(s(t)=F_Q(\bar x,B,t)-\dot{\bar x}\).  This form exposes the matrices
that a block-encoding, quantum linear-system solver, or QSP/QSVT matrix-function
routine would have to access.  A credible fault-tolerant claim must specify the
oracle model for \(J_Q\), the condition number or amplification factor, and the
observable extracted from the amplitude-encoded output.

\paragraph*{Fault-tolerant lifting layer.}
Because Eq.~\eqref{eq:paper_v_delta_rho12q} contains polynomial products of
the quantum-target variables, one can also consider a Carleman-type lift.  If
\begin{equation}
\begin{aligned}
F_Q(x,B,t)
&=F_0(B,t)+A_1(B,t)x
\\
&\quad
+A_2(B,t)x^{\otimes2}
+A_3(B,t)x^{\otimes3},
\end{aligned}
\end{equation}
then define
\begin{equation}
X_K=\left(x,\ x^{\otimes2},\ldots,x^{\otimes K}\right).
\end{equation}
The lifted system is approximately linear,
\begin{equation}
\dot X_K=\mathcal A_K(B,t)X_K+\tau_K ,
\end{equation}
where \(\tau_K\) is the truncation remainder.  This is the cleanest route from
the nonlinear hierarchy to fault-tolerant linear-ODE machinery, but its cost is
controlled by the lifted dimension, sparsity, dissipativity, truncation error,
and output model.  It should remain a theory route until those quantities are
bounded for the tested regimes.

\paragraph*{Decision rule.}
The working hierarchy is therefore:
\begin{enumerate}
\item characterize the classical \(L=2\) ODE and its physicality defects;
\item test NISQ only as a residual, tangent-projection, or observable-estimation
method for the \(x\)-sector with classical \(B(t)\);
\item analyze fault-tolerant routes only after linearization, implicit-step
linear solves, or controlled polynomial lifting;
\item avoid claiming a full quantum solution of the coupled system unless
Eq.~\eqref{eq:paper_v_B}, nonlinear-product access, and output extraction are
all included in the cost model.
\end{enumerate}

\begin{thebibliography}{9}

\bibitem{RivaSimoniPing2026}
G. Riva, J. Simoni, and Y. Ping,
``Open-quantum-system theory of non-Markovian electron--phonon dynamics,''
arXiv:2606.22233 (2026).

\bibitem{McLachlan1964}
A. D. McLachlan,
``A variational solution of the time-dependent Schrodinger equation,''
Mol. Phys. \textbf{8}, 39 (1964).

\bibitem{LiBenjamin2017}
Y. Li and S. C. Benjamin,
``Efficient variational quantum simulator incorporating active error minimization,''
Phys. Rev. X \textbf{7}, 021050 (2017).

\bibitem{Yuan2019}
X. Yuan, S. Endo, Q. Zhao, Y. Li, and S. C. Benjamin,
``Theory of variational quantum simulation,''
Quantum \textbf{3}, 191 (2019).

\end{thebibliography}

\end{document}
